\chapter*{Agradecimientos}
\addcontentsline{toc}{chapter}{Agradecimientos}

Los autores expresan su más profundo agradecimiento a la directora de este proyecto integrador y directora de LIDeSIA, la Dra.~Laura Cecilia Diaz Dávila, por su liderazgo técnico, su visión en el campo de la ingeniería y la inteligencia artificial, y por haberles confiado este desafiante proyecto. Al co-tutor, el Ing.~Ángel Nicolás Heredia, por su acompañamiento. A la Secretaria de Ciencia y Técnica de la Facultad de Ciencias Económicas (FCE), la Dra.~María Inés Stimolo, por su valioso aporte analítico y metodológico y su constante respaldo académico. La combinación de sus aportes y su compromiso con la excelencia fueron fundamentales para el desarrollo de esta investigación y para la formación profesional de los autores.

Este reconocimiento se extiende a los integrantes del Laboratorio de Investigación y Desarrollo de Software e Inteligencia Artificial y al personal del Ministerio Público Fiscal que participó del proyecto, por su constante acompañamiento durante este trabajo.

Este trabajo utilizó recursos computacionales provistos por UNC Supercómputo (CCAD), Universidad Nacional de Córdoba, Argentina, en el marco del Programa Sistema Nacional de Infraestructuras de Investigación Científico-Tecnológica (SNIICTs).
